\begin{lemma}[A finite shift orbit is decided by one finite prefix]
Let $h$ be locally constant, let $x$ be an increasing enumeration, and let
$m\in\mathbb N$.  There is a finite prefix of $x$ which determines all of
the finitely many values
\[
h(x),h(\mathsf{ShiftMap}(x)),\ldots,
h(\mathsf{PowerShiftIter}(m,x)).
\]
Precisely, there is $N$ such that whenever
$\mathsf{PrefixAgree}(N,x,y)$, then for every $j\leq m$,
\[
h(\mathsf{PowerShiftIter}(j,y))
=h(\mathsf{PowerShiftIter}(j,x)).
\]
\end{lemma}
\begin{proof}
For every $j\leq m$, apply the definition of local constancy
\noderef{LocallyConstant} at $\mathsf{PowerShiftIter}(j,x)$.  Thus, for each
such $j$, choose $n_j\in\mathbb N$ so that every $z$ satisfying
\[
\mathsf{PrefixAgree}(n_j,\mathsf{PowerShiftIter}(j,x),z)
\]
satisfies
\[
h(z)=h(\mathsf{PowerShiftIter}(j,x)).
\]
There are finitely many indices $j\leq m$, so these choices form the finite
set
\[
A=\{n_j+j\mid j\leq m\}\subseteq\mathbb N.
\]
The set $A$ is nonempty because $0\leq m$.  Let $N=\max A$.  By the defining
property of this maximum, for every $j\leq m$,
\[
n_j+j\leq N. \tag{1}
\]

Now suppose that $\mathsf{PrefixAgree}(N,x,y)$, and fix an arbitrary
$j\leq m$.  If $i<n_j+j$, then (1) gives $i<N$.  Hence, by the definition of
\noderef{PrefixAgree}, the assumed agreement at length $N$ restricts to
\[
\mathsf{PrefixAgree}(n_j+j,x,y). \tag{2}
\]
Apply \noderef{PowerShiftIterPrefixAgree} to (2), with prefix parameter
$n_j$ and shift parameter $j$.  This gives
\[
\mathsf{PrefixAgree}(n_j,\mathsf{PowerShiftIter}(j,x),
  \mathsf{PowerShiftIter}(j,y)). \tag{3}
\]
The choice of $n_j$ and (3) therefore give
\[
h(\mathsf{PowerShiftIter}(j,y))
 =h(\mathsf{PowerShiftIter}(j,x)).
\]
Because the same maximum $N$ was chosen before $j$ and the argument applies
to every $j\leq m$, this one $N$ works simultaneously for all the required
indices, proving the lemma.
\end{proof}
